% Created 2023-04-26 mié 17:55
% Intended LaTeX compiler: pdflatex
\documentclass[11pt]{article}
\usepackage[utf8]{inputenc}
\usepackage[T1]{fontenc}
\usepackage{graphicx}
\usepackage{longtable}
\usepackage{wrapfig}
\usepackage{rotating}
\usepackage[normalem]{ulem}
\usepackage{amsmath}
\usepackage{amssymb}
\usepackage{capt-of}
\usepackage{hyperref}
\usepackage{minted}

\renewcommand*{\contentsname}{Contenidos}
\renewcommand*{\figurename}{Figura}
\renewcommand*{\tablename}{Tabla}
\usepackage{parskip}
\usepackage[utf8]{inputenc} %% For unicode chars
\usepackage{placeins}
\usepackage[margin=2.50cm]{geometry}
\usepackage[T1]{fontenc}
\usepackage{mathpazo}
\linespread{1.05}
\usepackage[scaled]{helvet}
\usepackage{courier}
\usepackage[table]{xcolor}
\hypersetup{colorlinks=true,linkcolor=blue}
\RequirePackage{fancyvrb}
\DefineVerbatimEnvironment{verbatim}{Verbatim}{fontsize=\small,formatcom = {\color[rgb]{0.5,0,0}}}
\usepackage{mdframed}
\BeforeBeginEnvironment{minted}{\begin{mdframed}}
\AfterEndEnvironment{minted}{\end{mdframed}}
\setcounter{secnumdepth}{4}
\date{}
\title{Automatización de tareas}
\hypersetup{
 pdfauthor={},
 pdftitle={Automatización de tareas},
 pdfkeywords={},
 pdfsubject={},
 pdfcreator={Emacs 27.1 (Org mode 9.6)}, 
 pdflang={Spanish}}
\begin{document}

\maketitle
\setcounter{tocdepth}{4}
\tableofcontents

\setlength\parindent{10pt}



\section{Automatización de tareas}
\label{sec:orgf641fbc}

\subsection{Airflow}
\label{sec:orgccca450}
Airflow es una plataforma de orquestación de flujo de trabajo que se
utiliza para automatizar y programar tareas en un entorno de
datos. Permite crear flujos de trabajo (DAG) de manera fácil y rápida,
lo que puede ser muy útil en el contexto de machine learning.

\subsection{Microsoft Power Auomate}
\label{sec:org8054020}
\subsection{Luigi}
\label{sec:org6795d89}
\subsection{Dagster}
\label{sec:org1a4d3f1}


\section{Tarea}
\label{sec:org4c16067}
Un pequeño flujo de trabajo con Airflow relacionado con machine
learning podría involucrar los siguientes pasos:

\begin{enumerate}
\item Extracción y preprocesamiento de datos: Crear un DAG que extraiga
datos de una fuente externa, como una base de datos o un archivo
CSV. Una vez que se extraen los datos, se pueden realizar
operaciones de preprocesamiento para preparar los datos para el
modelado. Esto podría incluir la limpieza de datos faltantes, la
eliminación de variables innecesarias, la codificación de
variables categóricas, y la normalización de variables numéricas.
\item Entrenamiento del modelo: Crear una tarea en el DAG que entrene un
modelo de machine learning utilizando los datos
preprocesados. Esto puede implicar la división de los datos en
conjuntos de entrenamiento y prueba, la selección de un algoritmo
de machine learning adecuado y la optimización de sus parámetros.
\item Validación del modelo: Una vez que se ha entrenado el modelo,
crear una tarea que realice la validación del modelo con los datos
de prueba. Esto puede involucrar la evaluación del rendimiento del
modelo utilizando métricas de evaluación, como el error cuadrático
medio o la precisión.
\item Despliegue del modelo: Una vez que se ha validado el modelo, crear
una tarea que implemente el modelo en un entorno de
producción. Esto puede implicar el empaquetamiento del modelo en
un formato que pueda ser utilizado por otros sistemas, como una
API RESTful, o la integración del modelo en una plataforma de
nube, como AWS o Google Cloud.
\item Monitoreo del modelo: Finalmente, crear una tarea que monitoree
continuamente el rendimiento del modelo en producción, y alerte a
los usuarios si el rendimiento del modelo disminuye o si se
detectan problemas.
\end{enumerate}

\subsection{Ejemplo Airflow}
\label{sec:org3aa5928}

\begin{enumerate}
\item Crear el archivo DAG
Primero, crearemos un archivo Python que definirá nuestro DAG. Este
archivo puede llamarse, por ejemplo, iris\_ml\_workflow.py. A
continuación se muestra el código que debe contener:
\begin{minted}[]{python}
from datetime import datetime, timedelta
from airflow import DAG
from airflow.operators.python_operator import PythonOperator

# Importar las funciones que definiremos más adelante
from data_extraction import extract_data
from data_preprocessing import preprocess_data
from train_model import train_model
from validate_model import validate_model

default_args = {
    'owner': 'airflow',
    'depends_on_past': False,
    'start_date': datetime(2023, 4, 28),
    'email': ['airflow@example.com'],
    'email_on_failure': False,
    'email_on_retry': False,
    'retries': 1,
    'retry_delay': timedelta(minutes=5),
}

dag = DAG(
    'iris_ml_workflow',
    default_args=default_args,
    description='A workflow to train and validate'
    'a machine learning model for Iris classification',
    schedule_interval=timedelta(days=1),
)

# Definir las tareas del DAG
extract_task = PythonOperator(task_id='extract_data',
                       python_callable=extract_data, dag=dag)
preprocess_task = PythonOperator(task_id='preprocess_data',
                       python_callable=preprocess_data, dag=dag)
train_task = PythonOperator(task_id='train_model',
                       python_callable=train_model, dag=dag)
validate_task = PythonOperator(task_id='validate_model',
                       python_callable=validate_model, dag=dag)

# Establecer las dependencias entre las tareas
preprocess_task.set_upstream(extract_task)
train_task.set_upstream(preprocess_task)
validate_task.set_upstream(train_task)
\end{minted}

En este código, creamos un DAG con el nombre
"iris\_ml\_workflow". Definimos cuatro tareas: "extract\_data",
"preprocess\_data", "train\_model" y "validate\_model". También
establecemos las dependencias entre las tareas.
\end{enumerate}


\begin{enumerate}
\item Crear las funciones para las tareas
\end{enumerate}

A continuación, crearemos las funciones que se ejecutarán en cada
tarea. Suponemos que tienes los archivos data\_extraction.py,
data\_preprocessing.py, train\_model.py y validate\_model.py en el mismo
directorio que el archivo iris\_ml\_workflow.py. A continuación se
muestra el código para cada archivo.
data\_extraction.py:
\begin{minted}[]{python}
import pandas as pd
from sklearn.datasets import load_iris

def extract_data():
    iris = load_iris()
    data = pd.DataFrame(iris.data, columns=iris.feature_names)
    target = pd.DataFrame(iris.target, columns=['target'])
    data.to_csv('/path/to/data/iris_data.csv', index=False)
    target.to_csv('/path/to/data/iris_target.csv', index=False)

\end{minted}

data\_preprocessing.py:
\begin{minted}[]{python}
import pandas as pd
from sklearn.model_selection import train_test_split
from sklearn.preprocessing import StandardScaler

def preprocess_data():
    data = pd.read_csv('/path/to/data/iris_data.csv')
    target = pd.read_csv('/path/to/data/iris_target.csv')
    X_train, X_test, y_train, y_test =
             train_test_split(data, target, test_size=0.2)
    scaler = StandardScaler()
    X_train_scaled = scaler.fit_transform(X_train)
    X_test_scaled = scaler.transform(X_test)
    X_train_scaled.to_csv('/path/to/data/iris_train_data.csv', index=False)
    y_train.to_csv('/path/to/data/iris_train_target.csv', index=False)
    X_test_scaled.to_csv('/path/to/data/iris_test_data.csv', index=False)
    y_test.to_csv('/path/to/data/iris_test_target.csv', index=False)


\end{minted}
`train\_model.py`:
\begin{minted}[]{python}
import pandas as pd
from sklearn.linear_model import LogisticRegression

def train_model():
    X_train = pd.read_csv('/path/to/data/iris_train_data.csv')
    y_train = pd.read_csv('/path/to/data/iris_train_target.csv')
    clf = LogisticRegression()
    clf.fit(X_train, y_train)
    joblib.dump(clf, '/path/to/models/iris_model.joblib')
\end{minted}

validate\_model.py:
\begin{minted}[]{python}
import pandas as pd
from sklearn.metrics import accuracy_score
import joblib

def validate_model():
    X_test = pd.read_csv('/path/to/data/iris_test_data.csv')
    y_test = pd.read_csv('/path/to/data/iris_test_target.csv')
    clf = joblib.load('/path/to/models/iris_model.joblib')
    y_pred = clf.predict(X_test)
    accuracy = accuracy_score(y_test, y_pred)
    with open('/path/to/results/accuracy.txt', 'w') as f:
        f.write(str(accuracy))

\end{minted}
En estas funciones, extraemos los datos del dataset Iris, los
preprocesamos, entrenamos un modelo de regresión logística y validamos
el modelo en un conjunto de prueba. Los datos y modelos se guardan en
archivos en el sistema de archivos local.

3.Ejecutar el DAG Una vez que hayamos creado nuestro DAG y las
funciones para las tareas, podemos iniciar Airflow y ejecutar el
DAG. Para iniciar Airflow, ejecuta el siguiente comando en tu
terminal:

\begin{minted}[]{bash}
airflow webserver -p 8080
\end{minted}

Luego, en otra terminal, ejecuta el siguiente comando:

\begin{minted}[]{bash}
airflow scheduler
\end{minted}

Esto iniciará los servidores de Airflow y el planificador,
respectivamente. Luego, en tu navegador, ve a localhost:8080 para
acceder a la interfaz de usuario de Airflow. Allí, deberías ver el DAG
"iris\_ml\_workflow" en la pantalla. Haz clic en "Trigger DAG" para
iniciar el flujo de trabajo.

Si todo ha ido bien, deberías ver que cada tarea se ejecuta en
secuencia. Los archivos de datos, modelos y resultados deberían estar
disponibles en el sistema de archivos local, y deberías poder revisar
el archivo accuracy.txt para ver la precisión del modelo en el
conjunto de prueba.
\end{document}